% Limitations — Data in Brief fragment. Descriptive and honest; describes the
% data and its bounds only. The scientific RASAero comparison result is argued
% in the companion research article \cite{yu2026jsr}, NOT here.
% Every number matches paper/_shared/CANONICAL_FACTS.md and the published CSVs.

The following limitations bound the reuse of this dataset. They are stated
descriptively; the accuracy result they qualify is interpreted in the companion
research article~\cite{yu2026jsr}, which this data article accompanies and which
should be consulted for any scientific conclusion drawn from the tables.

\paragraph{Heterogeneous ground truth and an unsubtracted measurement floor.}
The measured apogees in \texttt{flight\_comparison.csv} are not produced by a
single instrument or campaign: the \texttt{flight\_data\_type} column records
barometric-altimeter, GPS, optical-tracking, and integrated-accelerometer
determinations across high-power-rocketry flights of differing instrumentation
and vintage, while the exploratory sounding-rocket apogees in
\texttt{sounding\_rockets\_exploratory.csv} are radar-, beacon-, and GPS-tracked
historical records. Each technique carries its own bias and dispersion, so the
ground-truth column embeds an irreducible measurement uncertainty of order
$1\%$. This floor is reported but deliberately \emph{not} subtracted from the
model errors: the published \texttt{err\_thiswork\_pct} and
\texttt{err\_rasaero\_pct} values are therefore conservative, in the sense that
some of the residual they attribute to the predictors is in fact ground-truth
noise. Users comparing any predictor against this corpus should treat
sub-percent differences as within the noise of the reference data.

\paragraph{Reconstruction uncertainty dominates the exploratory high-Mach set.}
The headline corpus (\texttt{flight\_comparison.csv}, $25$ flights, Mach
$0.54$--$4.33$) contains no flight above Mach $4.33$, and the high-Mach behavior
of the method is documented only by the separate, explicitly exploratory
\texttt{sounding\_rockets\_exploratory.csv} table. Of the $\sim$$20$ historical
sounding-rocket flights in that file, only three fall within $\pm10\%$ of the
recorded apogee --- Black Brant~V~VB (Mach $7.224$, $-6.97\%$), Nike--Deacon
no.~1 (Mach $4.956$, $-1.06\%$), and Nike--Deacon no.~2 (Mach $5.079$,
$-0.89\%$). The remainder scatter widely: the Nike--Apache family over-predicts
systematically ($+24$ to $+36\%$) and the Arcas and HEROS~3 vehicles
under-predict strongly ($-29$ to $-69\%$), with a small number of simulation-error
or zero-apogee cases reported transparently rather than removed. This spread is
driven far more by motor-thrust and geometry \emph{reconstruction} uncertainty in
the poorly documented historical source material than by the aerodynamic model:
these flights are provided so that the full, uncurated picture is on record, and
they must \emph{not} be read as a flight-validation result. Reporting only the
three admissible cases would constitute outcome-based selection; the entire set
is published for this reason.

\paragraph{RASAero II values are recorded predictions, not independent reruns.}
The \texttt{apogee\_rasaero\_ft} and \texttt{err\_rasaero\_pct} columns are the
predictions \emph{as published} in Rogers' public RASAero~II altitude-comparison
collection~\cite{rogers2015,rogers_rasaero_alt}; they are reproduced verbatim and
were \emph{not} regenerated by re-running RASAero~II on independently
reconstructed input decks. They are therefore a fixed, version-locked commercial
reference rather than a controlled head-to-head rerun. Any paired comparison a
user performs against the \texttt{apogee\_thiswork\_ft} column inherits this
asymmetry and should be framed accordingly.

\paragraph{The MESOS~293K model value is the corpus's largest single-flight error.}
The higher-Mach of the corpus's two two-stage vehicles (flight $25$, MESOS~293K,
the highest-Mach flight at Mach $4.33$; the other two-stage flight is the AeroPac~104K
Two-Stage at Mach $3.04$) is predicted by the archived code release at $273{,}056$~ft, a
signed error of $-6.96\%$---the largest single-flight error in the corpus. The
value is reproducible (confirmed by running the MESOS flight in isolation) and is
the value carried in the published CSV. An earlier deposited version (RFD~v1.2)
carried a different, erroneous value for this single flight; it has no defensible
derivation and is superseded by the version of record. The $-6.96\%$ value remains
within the $\pm10\%$ admission band, so the headline ``$25/25$ within $\pm10\%$''
count is unaffected.

\paragraph{Two base-drag constants make the companion-paper headline partly
in-sample.} Two base-drag scale constants used to generate the
\texttt{apogee\_thiswork\_ft} column (\texttt{THICK\_BL\_K} and
\texttt{SLENDER\_BODY\_K} in the drag calculator) are corpus-frozen: they were
calibrated against the apogee residuals of a subset of these very flights (Raven,
Rabia and its short-fin-can variant, Kinsel, and Torrent via the source
diagnostic). The aggregate accuracy of the \texttt{this work} column on the full
corpus is therefore \emph{partly in-sample}, and any statistic computed over all
$25$ flights should be read with that disclosure in mind. We do not re-argue the
generalization question here; the companion research article~\cite{yu2026jsr}
addresses it directly via a decontaminated prospective holdout (every flight any
constant touched moved into a development set, leaving a genuinely blind holdout).
Users who wish to evaluate the method out-of-sample should likewise exclude the
calibration flights listed above before scoring.
